%% ============================================================
%% 分片 04：原文第 3 节「Topological Field Theory From ${\EUN}=4$ Super Yang-Mills Theory」（原文 glshort.tex 起于第 1216 行）
%% 由 tools/harvmac2tex.py 从 arXiv 源码（hep-th/0604151v3）自动生成。
%% 只译英文散文；公式、编号、\label、\cite 一律不改。
%% ============================================================

\section{Topological Field Theory From ${\EUN}=4$ Super
Yang-Mills Theory} \label{sec:tqftfour}

Our next goal is to describe how to construct a four-dimensional
topological quantum field theory (TQFT) on any
four-manifold\footnote{We generally assume $M$ to be oriented, as the
construction is more interesting in this case.  However, if we
specialize to the case that the theta angle of eqn. \eqref{eq:oopy}\
vanishes and the canonical parameter $\Psi$ of section 第~\ref{sec:canonpar}~节\
equals 0 or $\infty$, then the construction makes sense even for
unorientable $M$.} $M$ by twisting of ${\EUN}=4$ super Yang-Mills
theory. In constructing TQFT's, it is most natural to use
Euclidean signature. Topological field theory is most naturally
related to Euclidean signature, and in any event the twisting we
use does not work well in Lorentz signature. When we specialize to
$M=\Bbb{R}\times W$ or $\Bbb{R}^2\times C$ for a three-manifold
$W$ or a two-manifold $C$, we can usefully return to Lorentz
signature by taking Lorentz signature on $\Bbb{R}$ or $\Bbb{R}^2$.

One important change in going to Euclidean signature is that the
spin representations $(\btwo,\bone)$ and $(\bone,\btwo)$ of
$Spin(4)$ are pseudoreal, while for $Spin(1,3)$ these
representations are complex conjugates of each other.
Correspondingly,
 the operator
$\Gamma_0\Gamma_1\Gamma_2\Gamma_3$ that distinguishes the two spin
representations now squares to $+1$; to avoid confusion, we will
call it $\hat\Gamma_E$ when using Euclidean signature and restrict
$\hat\Gamma$ for the Lorentz signature case.  Similarly, the two
spin representations ${\EUS}^\pm$ are complex conjugates of each
other in Euclidean signature, and the operator
$\Gamma_0\Gamma_1\dots\Gamma_9$ that distinguishes them, which we
now call $\bar\Gamma_E$, has eigenvalues $\mp i$ on ${\EUS}^\pm$.

To determine the sign, letting a subscript $L$ or $E$ refer to
Lorentz or Euclidean signature, we make the Wick rotation from
Lorentz to Euclidean signature by $ix^0_L=x^0_E$, so
$\partial/\partial x^0_L=i\partial/\partial x^0_E$ and
$\Gamma_{0\,L}=i\Gamma_{0\,E}$.  Hence, $\bar\Gamma_L =
i\bar\Gamma_E$, and as $\bar\Gamma=\bar\Gamma_L$ acts on
$\EUS^\pm$ as multiplication by $\pm 1$, $\bar\Gamma_E$ acts by
multiplication by $\mp i$.  In particular, if $\epsilon$ is a
supersymmetry generator, then
% 式 (3.1)
\setcounter{equation}{0}
\begin{equation}\label{eq:nohinks}
\bar\Gamma_E\epsilon=-i\epsilon.
\end{equation}
